\begin{frame}{$\tanh$를 거친 은닉 상태: 포화}
  실제로 $\tanh$를 거치면 $h_t = \tanh(z_t) =
  [0.197,\; 0.444,\; 0.462,\; 0.673]$.
  \begin{itemize}
    \item 이 $h_t$를 누적합 추정치로 ($\hat{s}_t = 10 \cdot h_t$):
      누적합 $[2, 5, 6, 10]$의 추정 $\to$
      $[1.97,\; 4.44,\; 4.62,\; 6.73]$
    \item 초기에는 거의 정확하고, 입력이 클수록 $\tanh$의 \textbf{포화}
      로 점점 아래로 벗어남 --- 두 원인:
    \begin{enumerate}
      \item $\tanh$의 ``요약 벡터는 $[-1, 1]$ 안에 갇힌다''는 한계
      \item 큰 숫자의 누적합을 \textbf{저차원 벡터 하나}에 담으려
        할 때의 왜곡
    \end{enumerate}
    \item 실전 RNN은 \textbf{여러 차원} 은닉 상태(11.2절 $H=8$) +
      11.3절의 \textbf{게이트} 구조로 이 문제를 다룬다
    \item \kq{주의: 11.2절의 ``abcabc…'' 모델은 이 누적합을 스스로
      학습한 것이 아니라 \textbf{단 3글자 주기라는 더 쉬운 요령}을
      배운다 --- 같은 구조라도 데이터에 따라 ``기억''의 내용이
      바뀐다.}
  \end{itemize}
\end{frame}
